% ============================================================================
% 第 4 章 - Verilog 芯片设计
% ============================================================================
\chapter{Verilog 芯片设计}

前面我们在用户态和裸机环境下写软件——我们一直假定 CPU、UART、RAM
这些东西已经存在。本章我们反过来：用 Verilog 写一个简单的硬件模块，
让自己理解 CPU 之外的世界是如何工作的。

我们写三个模块：
\begin{itemize}
    \item \textbf{adder.v}：一个 64 位字节加法器，演示如何用寄存器做简单计算
    \item \textbf{uart\_tx.v}：UART 发送器，把并行字节转成串行比特流送出
    \item \textbf{uart\_rx.v}：UART 接收器，把外部输入的串行比特流转回并行字节
\end{itemize}

\section{Verilog 基本概念}
% ============================================================================
% 4.1 Verilog 基本概念
% ============================================================================
\subsection{软件和硬件的根本区别}

汇编语言是给 CPU 看的——它描述的是"一步一步做什么"。
Verilog 是给综合器看的——它描述的是"硬件模块之间如何连接、每个时钟沿做什么"。

\paragraph{并行 vs 顺序}
在 C/汇编里，代码是顺序执行的。Verilog 不一样：
每个 \texttt{always} 块都在\textbf{同时}运行。CPU 内部有几十上百个
这样的并行块。

\paragraph{两个基本类型}
\begin{itemize}
    \item \textbf{reg}：寄存器，在时钟沿保存值
    \item \textbf{wire}：连线，只是把一个模块的输出连到另一个模块的输入
\end{itemize}

\subsection{同步逻辑：时钟驱动一切}

几乎所有现代数字设计都使用"同步时序逻辑"：

\begin{lstlisting}[style=verilog]
always @(posedge clk or negedge rst_n)
begin
    if (!rst_n)
        q <= 1'b0;        // 复位
    else
        q <= d;           // 每个时钟上升沿，把 d 的值写入 q
end
\end{lstlisting}

\begin{itemize}
    \item \texttt{posedge clk}：时钟上升沿
    \item \texttt{negedge rst\_n}：复位（低电平有效）
    \item \texttt{<=}：非阻塞赋值，在组合逻辑里用 \texttt{=}，在时序逻辑里用 \texttt{<=}
    \item \texttt{1'b0}：表示 1 位二进制 0
\end{itemize}

\subsection{UART 通信的基本约定}

\begin{center}
\begin{tabular}{l l}
波特率 & 每秒发送的比特数，常见值如 115200 \\
数据位 & 一个字节通常 8 位 \\
停止位 & 1 位或 2 位，标识字节结束 \\
起始位 & 1 位，先把电平拉低来通知"我要开始发了" \\
\hline
总位数 & 1 起始 + 8 数据 + 1 停止 = 10 位 / 字节 \\
\end{tabular}
\end{center}

UART 不附带时钟信号，收发双方必须事先约定好同一个波特率，
然后自己在内部用时钟计数来估计每个比特的中心位置。


\section{adder.v：64 位字节加法器}
% ============================================================================
% 4.2 adder.v
% ============================================================================
\subsection{设计目标}

把 64 位输入（8 个字节）的每个字节提取出来，求和，输出一个 8 位结果。
这是一个纯组合逻辑（不需要时钟）+ 一个寄存器输出的简单示例。

\subsection{源码与讲解}

\begin{lstlisting}[style=verilog,caption=Chip/src/adder.v]
module adder
#(
    parameter CLK_FRE = 50,     // 时钟频率(MHz)
    parameter BAUD_RATE = 115200 // 波特率（这里未使用，但保持与其他模块一致）
)
(
    input clk,                   // 时钟输入
    input rst_n,                 // 异步复位，低电平有效
    input [63:0] in_register,    // 64 位输入（8 个字节）
    output reg [7:0] out_register // 8 位输出（求和结果）
);

always @(posedge clk or negedge rst_n)
begin
    if (!rst_n)
    begin
        out_register <= 8'd0;
    end
    else
    begin
        out_register
            <= (in_register >> (8*0) & 8'b11111111)
             + (in_register >> (8*1) & 8'b11111111)
             + (in_register >> (8*2) & 8'b11111111)
             + (in_register >> (8*3) & 8'b11111111)
             + (in_register >> (8*4) & 8'b11111111)
             + (in_register >> (8*5) & 8'b11111111)
             + (in_register >> (8*6) & 8'b11111111)
             + (in_register >> (8*7) & 8'b11111111);
    end
end
endmodule
\end{lstlisting}

\paragraph{位提取与掩码}
\begin{itemize}
    \item \texttt{in\_register >> (8*0)} 不移位；\texttt{\& 8'hFF} 只保留最低字节
    \item \texttt{in\_register >> (8*1)} 右移 8 位，保留第二字节
    \item \ldots 直到 \texttt{8*7} = 第 8 个字节（最高位）
\end{itemize}

每个字节都是 8 位，8 个这样的数相加，最大可能值是
\( 8 \times 255 = 2040 \)，需要至少 11 位存。但我们的
\texttt{out\_register} 只有 8 位——结果会\textbf{溢出并截断}，
只保留低 8 位。这是故意的——它演示了"硬件里做的运算宽度是
你事先规定好的，超出部分就丢掉"的特点。

\paragraph{综合后会是什么样的硬件}
综合器会把这个设计变成一个由 7 个 8 位加法器串起来的电路
（加法树），最后一级的和被一个 8 位寄存器锁住。每个时钟沿，
寄存器都会重新计算一次——如果输入不变，输出也不变。

如果你想保留完整的 11 位结果，只需把 \texttt{out\_register}
改成 \texttt{[10:0]}（11 位）。其余代码无需改动。

\subsection{一个更简洁的写法}

其实 Verilog 支持用 \texttt{for} 循环写重复逻辑：

\begin{lstlisting}[style=verilog,caption=用 generate 写的简化版本]
// generate 块在综合时被展开成具体的门电路
generate
    for (i = 0; i < 8; i = i + 1)
    begin : gen_bytes
        // 编译时展开：8 个独立的字节提取逻辑
    end
endgenerate
\end{lstlisting}

但对我们的目标而言——理解硬件在做什么——
直接写 8 次反而更清晰，因为你能看到每个字节在做同一件事。


\section{uart\_tx.v：UART 发送器}
% ============================================================================
% 4.3 uart_tx.v
% ============================================================================
\subsection{设计目标}

接收一个并行的 8 位字节（\texttt{tx\_data}），在 \texttt{tx\_data\_valid = 1}
时捕获它，然后按 UART 协议逐位发送到 \texttt{tx\_pin} 上。

协议：

一帧 11 位（起始位 + 8 数据位 + 停止位 + 空闲）：

\begin{center}
\begin{tabular}{lllllllllll}
位 & 0 & 1--8 & 9 & 10 \\
作用 & 起始 & 数据 b0--b7 & 停止 & 空闲 \\
电平 & 低 & (数据) & 高 & 高 \\
\end{tabular}
\end{center}

每一位的持续时间由波特率决定。在 50 MHz 时钟、115200 波特率下，
一位持续 \( \frac{50,000,000}{115200} \approx 434 \) 个时钟周期。

\subsection{状态机设计}

发送器使用 3 个状态的简单状态机：

\begin{center}
\begin{tabular}{l l}
\textbf{S\_IDLE} & 空闲。等待 \texttt{tx\_data\_valid}。 \\
\textbf{S\_START} & 发送起始位（拉低 tx\_pin 一位的时间）。 \\
\textbf{S\_SEND\_BYTE} & 逐位发送 8 个数据位，低位在前。 \\
\textbf{S\_STOP} & 发送停止位（拉高 tx\_pin 一位的时间），完成后回到 IDLE。 \\
\end{tabular}
\end{center}

\subsection{源码与讲解}

\begin{lstlisting}[style=verilog,caption=Chip/src/uart\_tx.v]
module uart_tx
#(
    parameter CLK_FRE = 50,      // 时钟频率(MHz)
    parameter BAUD_RATE = 115200 // 波特率
)
(
    input clk,
    input rst_n,
    input [7:0] tx_data,         // 要发送的字节
    input tx_data_valid,         // 数据有效：一次握手的开始
    output reg tx_data_ready,    // 告诉发送方：我已准备好接收下一个
    output tx_pin                // 串行输出引脚
);

localparam CYCLE = CLK_FRE * 1000000 / BAUD_RATE;
// CYCLE = 434 (50 MHz / 115200)

// 状态编码
localparam S_IDLE      = 1;
localparam S_START     = 2;
localparam S_SEND_BYTE = 3;
localparam S_STOP      = 4;

reg [2:0]  state;
reg [15:0] cycle_cnt;    // 每一位内的时钟计数（最多 65535）
reg [2:0]  bit_cnt;      // 已经发送了第几位（0..7）
reg [7:0]  tx_data_latch; // 锁存住要发送的字节
reg        tx_reg;        // 输出寄存器

assign tx_pin = tx_reg;   // tx_reg 的值直接送到 tx 引脚
\end{lstlisting}

状态转移逻辑：

\begin{lstlisting}[style=verilog]
always @(posedge clk or negedge rst_n)
begin
    if (!rst_n)
        state <= S_IDLE;
    else
        case (state)
            S_IDLE:
                if (tx_data_valid)
                    state <= S_START;
                else
                    state <= S_IDLE;
            S_START:
                if (cycle_cnt == CYCLE - 1)
                    state <= S_SEND_BYTE;
                else
                    state <= S_START;
            S_SEND_BYTE:
                if (cycle_cnt == CYCLE - 1 && bit_cnt == 3'd7)
                    state <= S_STOP;
                else
                    state <= S_SEND_BYTE;
            S_STOP:
                if (cycle_cnt == CYCLE - 1)
                    state <= S_IDLE;
                else
                    state <= S_STOP;
        endcase
end
\end{lstlisting}

数据准备好标志和锁存：

\begin{lstlisting}[style=verilog]
always @(posedge clk or negedge rst_n)
begin
    if (!rst_n)
        tx_data_ready <= 1'b0;
    else if (state == S_IDLE)
        tx_data_ready <= !(tx_data_valid);  // 空闲时才能接收
    else if (state == S_STOP && cycle_cnt == CYCLE - 1)
        tx_data_ready <= 1'b1;
end

always @(posedge clk or negedge rst_n)
begin
    if (!rst_n)
        tx_data_latch <= 8'd0;
    else if (state == S_IDLE && tx_data_valid)
        tx_data_latch <= tx_data;    // 在进入 START 状态时锁存数据
end
\end{lstlisting}

计数器（计算当前位已发送了多少个时钟周期）：

\begin{lstlisting}[style=verilog]
always @(posedge clk or negedge rst_n)
begin
    if (!rst_n)
        cycle_cnt <= 16'd0;
    else if ((state == S_START || state == S_SEND_BYTE ||
              state == S_STOP) && cycle_cnt < CYCLE - 1)
        cycle_cnt <= cycle_cnt + 16'd1;
    else
        cycle_cnt <= 16'd0;
end

always @(posedge clk or negedge rst_n)
begin
    if (!rst_n)
        bit_cnt <= 3'd0;
    else if (state == S_SEND_BYTE)
        if (cycle_cnt == CYCLE - 1)
            bit_cnt <= bit_cnt + 3'd1;
        else
            bit_cnt <= bit_cnt;
    else
        bit_cnt <= 3'd0;
end
\end{lstlisting}

最后，真正驱动输出引脚的逻辑：

\begin{lstlisting}[style=verilog]
always @(posedge clk or negedge rst_n)
begin
    if (!rst_n)
        tx_reg <= 1'b1;   // UART 空闲时 tx_pin 为高电平
    else
        case (state)
            S_IDLE:      tx_reg <= 1'b1;   // 空闲=高
            S_START:    tx_reg <= 1'b0;   // 起始位=低
            S_SEND_BYTE:
                         tx_reg <= tx_data_latch[bit_cnt];  // 从 bit0 开始逐位发
            S_STOP:     tx_reg <= 1'b1;   // 停止位=高
            default:    tx_reg <= 1'b1;
        endcase
end
endmodule
\end{lstlisting}

\paragraph{锁存为什么重要}
\texttt{tx\_data\_latch} 在进入 S\_START 时把 \texttt{tx\_data} 的值存下来。
如果没有这一步，上游逻辑在发送过程中改变 \texttt{tx\_data}，
我们就会输出一个被污染的字节——先几位是旧值，后几位是新值。

\paragraph{为什么 UART 空闲时 tx\_pin 必须为高}
UART 协议里，"线路空闲"就是高电平。起始位是把线路拉低一位时间，
这样接收端才能检测到"有数据开始到来"。如果空闲时电平是低，
接收端会把空闲状态误当做起始位，整个通信就错乱了。


\section{uart\_rx.v：UART 接收器}
% ============================================================================
% 4.4 uart_rx.v
% ============================================================================
\subsection{设计目标}

与发送器相反：持续监视 \texttt{rx\_pin}。当 \texttt{rx\_pin} 从 1 变 0
（起始位到来）时，启动一个计数器。在每个比特的中心位置采样
一次电平，攒够 8 位后输出 \texttt{rx\_data} 并拉高 \texttt{rx\_data\_valid}。

接收器比发送器难，因为：
\begin{itemize}
    \item 发送器是自己决定什么时候发一位，时钟完全可控
    \item 接收器不知道对方何时开头发送，必须自己检测起始位
    \item 接收器不能保证自己的时钟和对方时钟完全对齐，
        必须在比特\textbf{中心}采样以获得最大噪声容限
\end{itemize}

\subsection{状态机设计}

\begin{center}
\begin{tabular}{l l}
\textbf{S\_IDLE} & 监视 rx\_pin。检测到下降沿 → 进入 START 状态。 \\
\textbf{S\_START} & 计数半个周期，再采样一次确认真的是低电平（防毛刺）。 \\
\textbf{S\_REC\_BYTE} & 每过一个周期采样一位，共 8 位，低位在前。 \\
\textbf{S\_STOP} & 等待停止位，完成后回到 IDLE。 \\
\textbf{S\_DATA} & 拉高 valid 标志，让上层模块取走字节。 \\
\end{tabular}
\end{center}

\subsection{源码与关键思路讲解}

\begin{lstlisting}[style=verilog,caption=Chip/src/uart\_rx.v]
module uart_rx
#(
    parameter CLK_FRE = 50,
    parameter BAUD_RATE = 115200
)
(
    input clk,
    input rst_n,
    output reg [7:0] rx_data,       // 收到的字节
    output reg rx_data_valid,       // 一次有效，表明有新字节
    input rx_data_ready,             // 上层已取走，让我进入下一轮
    input rx_pin                     // 串行输入引脚
);

localparam CYCLE = CLK_FRE * 1000000 / BAUD_RATE;
localparam S_IDLE      = 1;
localparam S_START     = 2;
localparam S_REC_BYTE  = 3;
localparam S_STOP      = 4;
localparam S_DATA      = 5;

reg [2:0] state;
reg [2:0] next_state;
reg       rx_d0;           // 打一拍，同步 rx_pin
reg       rx_d1;           // 再打一拍，用于检测下降沿
wire      rx_negedge;      // 检测到一个下降沿
reg [7:0] rx_bits;         // 保存逐位收到的数据
reg [15:0] cycle_cnt;      // 当前位已持续了多久
reg [2:0] bit_cnt;         // 已收到第几位
\end{lstlisting}

\paragraph{打两拍：跨时钟域的基本操作}
\texttt{rx\_pin} 是一个从外部引脚进入的信号，它的变化不一定对齐
我们的时钟沿。直接用它来驱动逻辑会有"亚稳态"风险。
解决办法是\textbf{打两拍}：

\begin{lstlisting}[style=verilog]
assign rx_negedge = rx_d1 && ~rx_d0;  // 检测到从 1→0

always @(posedge clk or negedge rst_n)
begin
    if (!rst_n)
    begin
        rx_d0 <= 1'b0;
        rx_d1 <= 1'b0;
    end
    else
    begin
        rx_d0 <= rx_pin;
        rx_d1 <= rx_d0;
    end
end
\end{lstlisting}

打两拍后，我们内部观察到的信号变化至少延迟了两个时钟周期，
但它是\textbf{稳定的}——不会在一个时钟沿前后出现混乱的值。

\paragraph{状态转移逻辑：使用两个 always 块}
接收器用了"状态 / 次态分离"的写法，这是更标准的状态机写法：

\begin{lstlisting}[style=verilog]
// 第一块：状态更新（时序逻辑，只在 clk 沿更新）
always @(posedge clk or negedge rst_n)
begin
    if (!rst_n)
        state <= S_IDLE;
    else
        state <= next_state;
end

// 第二块：计算 next_state（组合逻辑，异步计算）
always @(*)
begin
    case (state)
        S_IDLE:
            if (rx_negedge)
                next_state = S_START;
            else
                next_state = S_IDLE;
        S_START:
            if (cycle_cnt == CYCLE - 1)
                next_state = S_REC_BYTE;
            else
                next_state = S_START;
        S_REC_BYTE:
            if (cycle_cnt == CYCLE - 1 && bit_cnt == 3'd7)
                next_state = S_STOP;
            else
                next_state = S_REC_BYTE;
        S_STOP:
            if (cycle_cnt == CYCLE/2 - 1)  // 半位就进入数据锁存
                next_state = S_DATA;
            else
                next_state = S_STOP;
        S_DATA:
            if (rx_data_ready)
                next_state = S_IDLE;
            else
                next_state = S_DATA;
        default:
            next_state = S_IDLE;
    endcase
end
\end{lstlisting}

\paragraph{为什么在 S\_STOP 里只等半位}
我们是在起始位的\textbf{中间}检测到稳定低电平的。
后续每位也都是从位中心开始的一个完整周期采样。
但停止位其实不需要等完整的一位——只要确认它是高电平就完成了。
所以用 \texttt{CYCLE/2 - 1}，只等半个周期，
省出的时间让我们更快回到 IDLE 状态，不漏掉下一个字节。

\paragraph{数据与有效标志}

\begin{lstlisting}[style=verilog]
always @(posedge clk or negedge rst_n)
begin
    if (!rst_n)
        rx_data_valid <= 1'b0;
    else if (state == S_STOP && next_state != state)
        rx_data_valid <= 1'b1;
    else if (state == S_DATA && rx_data_ready)
        rx_data_valid <= 1'b0;
end

always @(posedge clk or negedge rst_n)
begin
    if (!rst_n)
        rx_data <= 8'd0;
    else if (state == S_STOP && next_state != state)
        rx_data <= rx_bits;  // 把 8 位并行输出
end
\end{lstlisting}

\paragraph{逐位接收}

\begin{lstlisting}[style=verilog]
always @(posedge clk or negedge rst_n)
begin
    if (!rst_n)
        cycle_cnt <= 16'd0;
    else if ((state == S_START || state == S_REC_BYTE ||
              state == S_STOP) && cycle_cnt < CYCLE - 1)
        cycle_cnt <= cycle_cnt + 16'd1;
    else
        cycle_cnt <= 16'd0;
end

always @(posedge clk or negedge rst_n)
begin
    if (!rst_n)
        rx_bits <= 8'd0;
    else if (state == S_REC_BYTE && cycle_cnt == CYCLE/2 - 1)
        rx_bits[bit_cnt] <= rx_pin;   // 在每个比特的中心采样
end

endmodule
\end{lstlisting}

\paragraph{中心采样的重要性}
在 \texttt{cycle\_cnt == CYCLE/2 - 1} 时采样，相当于在
每个位的正中间读取电平。UART 通信两端的时钟总有微小的误差，
但只要误差不累计到超过半个比特周期，中心采样就可以正确读取数据。
这是 UART 协议的核心鲁棒性来源。

如果在边沿采样，发送器和接收器的时钟漂移累积起来，几个字节之后
你就会开始读错——因为你采样的位置在两个比特的交界处，
读出的电平取决于先到哪一侧。


\section{综合到 FPGA：从代码到硬件}
% ============================================================================
% 4.5 综合到 FPGA
% ============================================================================
\subsection{从 Verilog 到真实硬件的流程}

写好的 Verilog 代码通过下面几个步骤变成可以烧到 FPGA 里的文件：

\begin{enumerate}
    \item \textbf{综合（Synthesis）}：把 Verilog 代码翻译成一个由逻辑门、
        触发器、查找表（LUT）组成的网表
    \item \textbf{布局（Placement）}：决定每个逻辑门放在 FPGA 的哪个位置
    \item \textbf{布线（Routing）}：在 FPGA 的可编程互连资源里
        把需要通信的逻辑门连接起来
    \item \textbf{时序分析（Timing Analysis）}：验证设计能在目标时钟频率下
        正确工作——每条路径的传播延迟必须小于一个时钟周期
    \item \textbf{生成比特流（Bitstream Generation）}：把最终配置信息
        打包成 FPGA 芯片可以理解的二进制文件
\end{enumerate}

你手工用的工具通常是 Vivado（Xilinx 系列 FPGA）、Quartus（Intel/Altera）
或者开源工具链（Yosys + nextpnr + icestorm，针对 Lattice FPGA）。

\subsection{u\_lite 程序如何加载到 FPGA 的指令存储器}

在设计一个自制 CPU 时，指令内存有两种主要的初始化方式：

\paragraph{方式 1：利用 Block RAM 初始化}
在综合时，把你汇编产出的机器码写成一个初始化文件（.coe 或 .mem），
让综合器在编译时把这些值写进 FPGA 的 BRAM 里。

\begin{lstlisting}[style=verilog,caption=BRAM 里嵌入机器码的模板]
// 这个模块在综合时被工具识别为 "Single Port RAM with initial contents"
reg [31:0] instr_mem [0:1023];  // 1024 条 32 位指令 = 4 KB

initial begin
    $readmemh("code.mem", instr_mem);  // 综合时被解析
end

// 每个时钟沿读出一条指令
always @(posedge clk)
begin
    instr_out <= instr_mem[pc[11:2]];
end
\end{lstlisting}

在编译（汇编）层面，你把 \texttt{start.s} 这类文件用
\texttt{clang --target=aarch64-none-elf -c start.s -o start.o} 产出 .o，
再用 \texttt{llvm-objcopy -O binary} 产出纯二进制，最后
写一段脚本把它转成一行一个十六进制数的 .mem 文件。

\paragraph{方式 2：通过 UART bootloader}
更灵活的做法：在 FPGA 设计里放一个小的 bootloader（固定在 BRAM 里），
它在上电后立即启动，通过 UART 接收外部发来的程序二进制，
把它写入一块更大的 RAM/ROM，然后跳过去执行。

这对应 QEMU 章节里你所看到的思路：程序可以在运行时被写入硬件。
UART 发送器（uart\_tx）把程序字节送来，UART 接收器（uart\_rx）
把它们读出来，写进指令存储器——你就不需要重新综合了。

\subsection{和 QEMU 模拟的对应关系}

现在把第 3 章的 QEMU 裸机程序和本章的 Verilog UART 做个对比：

\begin{center}
\begin{tabular}{l l l}
\hline
层面 & QEMU 里 & FPGA 里对应的 \\
\hline
UART 设备 & 虚拟 PL011 & 你写的 UART 模块 \\
地址 & 内存路由到 PL011 & 地址解码器路由到 UART \\
CPU 取指 & TCG 翻译执行 ARM & 你未来写的 CPU 核 \\
DRAM & QEMU 模拟的内存 & FPGA 片上 BRAM 或 DDR \\
\hline
\end{tabular}
\end{center}

你现在写的 UART 模块，加上一个简单的 CPU，就是一个最小的 SoC。
把之前写的 QEMU 汇编程序（读取输入）改写成这个自制 CPU 可以
执行的指令格式后，就能在 FPGA 里跑起来。

\subsection{为什么要做这种事}

从学习角度，写过 UART 之后再回头看 QEMU 里的那段汇编：

\begin{lstlisting}[style=asm]
wait_uart:
    ldr  w3, [x0, #0x18]
    tbnz w3, #5, wait_uart
    str  w2, [x0]
\end{lstlisting}

你会立刻理解它的意思——这是在读取 FLAG 寄存器、检查发送满、
然后写数据寄存器。这些动作不是某个抽象的"系统调用"，
而是真正把字节送到一个你自己用 Verilog 写过的硬件模块里。
你知道电路里发生了什么，就知道软件层为什么要这么写。

反过来说，真正写过软件之后再去写硬件，你也清楚 UART 的数据寄存器
和状态寄存器为什么要分开成两个地址——因为软件会分开读它们。
两者是一体两面。

